\section{Semilinear descent and an adic bound}\label{sec:adic}
At height $a=p$ for $p\geq5$, and at height $a=p+1$ for $p\geq7$,
the upper arguments need continuous cohomology to vanish above degree
$a^2+1$, for both sphere coefficients
and finite-test coefficients. We first bound the cohomology of finite
adic quotients and then pass to their countable inverse limit. At height
$p$ the Galois quotient has order $p$, and its semilinear action on Witt
coefficients supplies the required descent. At height $p+1$ the full
extended group has finite $p$-cohomological dimension.

Put $k_a=\F_{p^a}$ and write
\[
 (E_a)_*=O_a[u^{\pm1}],\quad
 O_a=W(k_a)[[u_1,\ldots,u_{a-1}]],\quad
 \m_a=(p,u_1,\ldots,u_{a-1}),\quad |u|=-2,
\]
and $\G_a=\St_a\rtimes\Gal(k_a/\F_p)$. Here $\St_a$ is the
pure stabilizer over $k_a$.

Exactness of Galois invariants for $p$-complete twisted Witt modules is
proved in \cite[Lemma~1.14]{goerss-hopkins2020}. The explicit contraction
below applies to arbitrary semilinear modules.

\begin{lemma}[Trace-one contraction]\label{lem:semilinear}
Let a finite group $Q$ act on a commutative ring $W$. Suppose
$w\in W$ satisfies $\sum_{g\in Q}g(w)=1$. For any semilinear
$W$-module $P$, one has $H^i(Q,P)=0$ for $i>0$, with no finiteness
or flatness hypothesis on $P$.
In particular this holds for $W=W(\F_{p^a})$ and
$Q=\Gal(\F_{p^a}/\F_p)$, without a prime-to-$p$ hypothesis on $a$.
\end{lemma}
\begin{proof}
Use homogeneous group cochains: a degree-$q$ cochain is an equivariant
function $f:Q^{q+1}\to P$, with differential the alternating deletion
of coordinates. For $q\geq1$ define
\[
 (hf)(g_0,\ldots,g_{q-1})
   =\sum_{\gamma\in Q}\gamma(w)f(\gamma,g_0,\ldots,g_{q-1}).
\]
For $g\in Q$, replacing $\gamma$ by $g\gamma$ proves that $hf$ is equivariant:
semilinearity converts $(g\gamma)(w)\,g f$ into
$g(\gamma(w)f)$. In $hd+dh$, the terms deleting one of the displayed
$g_i$ cancel in pairs. The term deleting the inserted coordinate is
$\sum_\gamma\gamma(w)f=f$. Thus $hd+dh=\mathrm{id}$ in positive
degrees, proving acyclicity for arbitrary $P$.

For the Witt-ring case, the finite-field trace is onto. Lift a residue
trace-one element to $w_0\in W$. Its Witt trace lies in $\Z_p$ and is
congruent to one modulo $p$, hence is a unit. Dividing $w_0$ by this
invariant unit supplies $w$. We have not divided by $|Q|$.
\end{proof}

\begin{lemma}[Finite quotients and the extra adic degree]\label{lem:adic}
Let $G=S_0\rtimes Q$ be profinite with $Q$ finite, and let $D\geq0$
be an integer such that $\cd_p(S_0)\leq D$ for discrete $p$-primary modules.
Let $\{M_j\}_j$ be a countable tower of finite discrete $p$-primary
$G$-modules with surjective transitions, and give
$M=\varprojlim_j M_j$ the inverse-limit topology. Suppose each $M_j$ is a
$W$-module, $S_0$ acts $W$-linearly, $Q$ acts semilinearly, and $W$ admits the
trace-one element of Lemma~\ref{lem:semilinear}. Then
\[
 H^s_{\cts}(G,M_j)=0\quad(s>D),\qquad
 H^s_{\cts}(G,M)=0\quad(s>D+1).
\]
The second conclusion also holds if the first displayed finite-quotient
vanishing is supplied directly, without a semidirect-product hypothesis.
\end{lemma}
\begin{proof}
The pure-subgroup cochain complex is $W$-linear and has semilinear
$Q$-action, including the conjugation action on its arguments: for an
inhomogeneous $b$-cochain,
\[
 (\gamma f)(s_1,\ldots,s_b)=
 \gamma\bigl(f(\gamma^{-1}s_1\gamma,\ldots,
                  \gamma^{-1}s_b\gamma)\bigr).
\]
The pure differential is $W$-linear and commutes with this action. Its
cohomology is therefore semilinear, whether or not that cohomology has
separately been shown finite. The continuous Hochschild--Serre sequence
at the discrete coefficient level is
\[
 H^i\bigl(Q,H^b_{\cts}(S_0,M_j)\bigr)
 \Longrightarrow H^{i+b}_{\cts}(G,M_j).
\]
Lemma~\ref{lem:semilinear} kills its positive $Q$-columns, giving
\[
 H^s_{\cts}(G,M_j)=H^s_{\cts}(S_0,M_j)^Q.
\]
This proves the first bound. Continuous cochains into $M$ are compatible
cochains into its finite quotients. Their transition maps are degreewise
surjective: a continuous cochain has a finite value set, and any choice
of lifts of those values gives a continuous cochain into the next finite
quotient. No equivariant section is needed for inhomogeneous cochains.
The countable inverse-limit sequence is therefore
\begin{equation}\label{eq:adic-milnor}
 0\longrightarrow\varprojlim_j{}^1 H^{s-1}_{\cts}(G,M_j)
 \longrightarrow H^s_{\cts}(G,M)
 \longrightarrow\varprojlim_j H^s_{\cts}(G,M_j)\longrightarrow0.
\end{equation}
Both outer terms vanish for $s>D+1$. At the remaining boundary degree,
the same sequence gives
\[
 H^{D+1}_{\cts}(G,M)
 \cong\varprojlim_j{}^1 H^D_{\cts}(G,M_j).
\]
Thus the bound retains the possible extra inverse-limit degree.
\end{proof}

\begin{proposition}[The two upper heights]\label{prop:cohomology}
Let $V_0$ be a finite ordinary $p$-local spectrum, let $DV_0=F(V_0,S)$,
and give $M=\pi_q(E_a\wedge DV_0)$ its $\m_a$-adic topology and
Morava action. Then
\[
 H^s_{\cts}(\G_a,M)=0\qquad(s>a^2+1)
\]
in each of the following cases: $a=p$ with $p\geq5$, and $a=p+1$
with $p\geq7$.
\end{proposition}
\begin{proof}
The pure stabilizer is the maximal-order unit group in the central
division algebra over $\Q_p$ of dimension $a^2$, so it is compact
$p$-adic analytic of that dimension \cite[Sections~2.1--2.2]{morava1985}.
It has no element $g$ of order $p$ in either stated range. Indeed,
$(g-1)\Phi_p(g)=0$ in a division algebra and $g\ne1$ imply
$\Phi_p(g)=0$. The polynomial $\Phi_p(X+1)$ is Eisenstein, so $g$
generates the degree-$(p-1)$ subfield $\Q_p(\zeta_p)$. The division
algebra would then be a vector space over this subfield, forcing
$p-1\mid a^2$. For $a=p$ this forces $p-1\mid1$; for $a=p+1$ it
forces $p-1\mid4$. These are impossible in the stated ranges.
The Lazard--Serre theorem therefore gives
$\cd_p(\St_a)=a^2$ for discrete $p$-primary coefficients
\cite[Section~1, Corollary~(1)]{serre1965}.

For each $q$, the coefficient module $M$ is finitely generated over the
complete Noetherian ring $O_a$, by the finite construction of $DV_0$.
It is complete and separated; its quotients $M_j=M/\m_a^jM$ are finite
discrete $p$-primary modules with surjective transitions. This uses no
flatness of $M$ and allows torsion and both parities. The maximal ideal
is invariant. For $V_0=S$, continuity of the action on these quotients
is part of the Morava action used by Devinatz--Hopkins
\cite[Theorem~1(ii) and Remark~1.3]{dh2004}. For every finite $V_0$,
the actual action and its continuity are also obtained by the ordinary
cooperation construction of Proposition~\ref{wup:cochains} below.

The pure stabilizer fixes $W(k_a)$. It fixes the residue field and thus
the unique Hensel lifts of the roots of $X^{p^a}-X$, which are the
Teichm\"uller representatives; continuity then fixes their convergent
Witt expansions. Hence its action on $M_j$ is Witt-linear, while
$Q=\Gal(k_a/\F_p)$ acts semilinearly by the Witt lift. At $a=p$,
Lemma~\ref{lem:adic} with $D=p^2$ gives the bound. The full group has
the split Galois subgroup of order $p$; this conclusion is specific to
the stated coefficients. The cohomological-dimension theorem was
applied only to $\St_p$.

At $a=p+1$ with $p\geq7$, the Galois quotient has order $p+1$, prime to $p$.
An order-$p$ element of $\G_a$ would therefore lie in $\St_a$, where
we have just excluded it. The full group is a finite extension of the
compact $p$-adic analytic pure group and has the same dimension $a^2$.
The same dimension theorem now applies to the full group itself and
gives $H^s_{\cts}(\G_a,M_j)=0$ for $s>a^2$. Apply only the
inverse-limit part of Lemma~\ref{lem:adic} to this finite-quotient
bound. In both cases the finite quotients vanish for $s>a^2$, and
\eqref{eq:adic-milnor} gives the asserted bound $s>a^2+1$ for $M$.
\end{proof}

For the sphere-kernel argument, take $V_0=S$, so that $M=(E_a)_q$.
The exponent argument uses the same bound for every finite $V_0$,
including modules with torsion and both parities. In both applications
the topology is the full maximal-ideal topology, as in the
continuous-cochain convention of Devinatz--Hopkins
\cite[Remark~1.3]{dh2004}.

\begin{remark}[A sharper line for sphere coefficients]\label{rem:sphere-sharp}
For the sphere homotopy fixed point spectral sequence, Bobkova and
collaborators state an $E_2$ vanishing line at $s=n^2+1$ when
$p-1\nmid n$ \cite[Equation~(2) and Remark~1.2]{bobkova2025}.
At $p=n=5$ this gives
$H^s_{\cts}(\G_5;(E_5)_q)=0$ for $s>25$.
That statement concerns the ordinary sphere coefficients $(E_n)_q$.
It does not assert the same line for every finite-test module
$\pi_q(E_n\wedge DV_0)$ or identify an actual image filtration.
Proposition~\ref{prop:cohomology} supplies the common conservative bound
used in this paper.
\end{remark}
